\subsection{Positive-current subtraction and divisorial residuals}\label{sec:nef-divisorial-residuals}

Another way to change the order of degeneration is to subtract a positive current from the background. If $\theta-\beta$ differs from a positive current by a $\ddc$-exact term, then every Gauduchon pairing of $\theta$ dominates the corresponding pairing of $\beta$. When the residual $\beta$ is still nef, its perturbed volumes are therefore smaller, and its orders are at least those of $\theta$. The main example is the subtraction of the curvature form of an effective divisor. Such a subtraction can preserve the order, raise it by an integer or by a fraction, or produce unequal lower and upper orders; and the residual need not be nef.

\Needspace{15\baselineskip}
\begin{proposition}\label{prop:nef-positive-current-subtraction}\index{Nef degeneration!positive-current subtraction}
Let $\theta,\beta$ be smooth nef real $(1,1)$-forms on a compact complex manifold $X$ of dimension $n$, and suppose that
\[
 S=(\theta-\beta)+\ddc\varphi\geq0
\]
is a positive $(1,1)$-current for some integrable function $\varphi$. For every Hermitian form $h$, every smooth strictly positive probability form $\mu$ and every $\epsilon>0$,
\begin{equation}\label{eq:nef-residual-volume-comparison}
 \mathcal V_\mu(\beta+\epsilon h)\leq \mathcal V_\mu(\theta+\epsilon h).
\end{equation}
Consequently $\underline a(\beta)\geq\underline a(\theta)$ and $\overline a(\beta)\geq\overline a(\theta)$.
\end{proposition}

\begin{proof}
For every Gauduchon metric $\gamma$, distributional integration by parts gives
\[
 \int_X(\theta-\beta)\wedge\gamma^{n-1}
 =\langle S,\gamma^{n-1}\rangle\geq0.
\]
Hence, for each test metric, the quotient for $\theta+\epsilon h$ is at least the quotient for $\beta+\epsilon h$, and taking infima gives $\lambda_\mu(\beta+\epsilon h)\leq\lambda_\mu(\theta+\epsilon h)$. Both root volumes are nonnegative, so raising to the $n$th power gives \eqref{eq:nef-residual-volume-comparison}. Dividing the logarithms by the negative number $\log\epsilon$ reverses the inequality and proves both comparisons of orders.
\end{proof}

For an effective real divisor $D$ with curvature form $\eta_D$, the Poincar\'e--Lelong formula gives $[D]=\eta_D+\ddc\varphi$ for a quasi-plurisubharmonic function $\varphi$. Thus \cref{prop:nef-positive-current-subtraction} applies to $\beta=\theta-\eta_D$, with $S=[D]$, whenever $\beta$ is nef. If $T=\theta+\ddc u$ is a positive current with $T-S\geq0$, then $T-S=\beta+\ddc(u-\varphi)$, so $\beta$ is the background of the positive current that remains after $S$ is removed from $T$.

\begin{proposition}\label{prop:nef-specified-residual-examples}
Subtracting the curvature form of an effective divisor from a smooth nef background, with a smooth nef residual, can preserve the common order, raise it by any integer between $1$ and the dimension or by a fraction, or produce a residual whose lower and upper orders differ. The residual can also fail to be nef.
\end{proposition}

\begin{proof}
\emph{Integer jumps.} On $E^n$, where $E$ is an elliptic curve, let $\omega_i$ be the pullback of a flat area form of integral one from the $i$th factor, and set $h=\sum_i\omega_i$ and $\mu=\omega_1\wedge\cdots\wedge\omega_n$. The divisor $F_i$ obtained by fixing a point in the $i$th factor has $\omega_i$ as a smooth curvature form. Indeed, on that elliptic curve any curvature form of the corresponding line bundle differs from $\omega_i$ by $\ddc$ of a solution of a zero-mean Poisson equation, so the metric can be adjusted on the curve and then pulled back. For $\theta=h$ and $D=\sum_{i=1}^kF_i$, the residual $\beta=\sum_{i=k+1}^n\omega_i$ is smooth and nef. The perturbed forms are closed, positive and of constant density with respect to $\mu$, so the zero potential solves their equations, and \cref{thm:gaud-volume-dual-formula} gives
\begin{equation}\label{eq:nef-integer-residual-jump}
 \mathcal V_\mu(\theta+\epsilon h)=n!(1+\epsilon)^n,\qquad
 \mathcal V_\mu(\beta+\epsilon h)=n!\epsilon^k(1+\epsilon)^{n-k}.
\end{equation}
These examples realize each integer jump $1,\ldots,n$. A nonzero divisor with jump zero is given by $\theta=2h$ and $D=F_1$, since then the residual is still strictly positive.

\emph{A fractional jump.} Use the background $\beta=b\omega_1+q\omega_2$ of \cref{prop:nef-diagonal-fractional-and-zero}, where $q=|\zeta|^{2m}$ near its unique zero and $m>1$. Choose $c>0$ and put $\theta=\beta+c\omega_2>0$ and $D=c(E\times\{w_0\})$, whose smooth curvature form is $c\omega_2$. The original background is strictly positive, so both its orders are zero, whereas the residual has common order $1-1/m$.

\emph{Unequal orders.} Let $\beta_{\mathrm{res}}$ be the oscillatory rank-one background of \cref{prop:nef-oscillatory-orders}, and put
\[
 \theta=\beta_{\mathrm{res}}+ch>0,\qquad
 D=c(\{z_0\}\times E)+c(E\times\{w_0\}),\qquad
 \eta_D=ch.
\]
The background $\theta$ has both orders zero: since $\theta\geq ch$, every test quotient is bounded below by a positive constant, while one fixed test gives a finite upper bound. Its smooth nef residual is exactly $\beta_{\mathrm{res}}$, whose lower and upper orders differ. In each of these examples, the residual plus $[D]$ is a positive current in the smooth $\ddc$-class of $\theta$, and its divisorial part is exactly $[D]$ because the residual is smooth. We do not need this current to have minimal singularities; that case is treated in \cref{cor:nef-boundary-minimal-removal}.

\emph{A residual that is not nef.} Blow up $\mathbb P^2$ at two distinct points. Let $H$ be the pullback of the class of the normalized Fubini--Study form, let $E_1,E_2$ be the exceptional curves, and let $L$ be the strict transform of the line through the two points. The pullback $\theta$ of the Fubini--Study form is smooth, semipositive and nef. Subtract curvature forms of $D=E_1+E_2$. The residual $\beta$ has class $H-E_1-E_2$, which is the class of $L$, and the Poincar\'e--Lelong formula for the total transform of the line shows that it contains the positive current $[L]$. The intersection numbers $H^2=1$, $H\cdot E_i=0$, $E_i^2=-1$ and $E_1\cdot E_2=0$ give
\[
 \int_L\beta=(H-E_1-E_2)^2=-1.
\]
If $\beta$ were nef, then, for a Hermitian form $h$, integrating smooth positive representatives $\beta+\epsilon h+\ddc u_\epsilon$ over $L$ would give $-1+\epsilon\int_Lh\geq0$ for all $\epsilon>0$, which is impossible.
\end{proof}

For a prime divisor $D$ and $\delta>0$, consider the infimum of the generic Lelong numbers along $D$ of the positive currents $\theta+\delta h+\ddc u$; its limit as $\delta\downarrow0$ is the minimal multiplicity of $\theta$ along $D$. For a nef background all minimal multiplicities vanish, since every $\theta+\delta h$ has a smooth positive representative. Removing the divisor defined by these multiplicities therefore leaves a nef background unchanged. Removing the divisorial part of a particular positive current in the class of $\theta$ is a different operation, even when that current has minimal singularities.

\begin{corollary}\label{cor:nef-boundary-minimal-removal}
On the ruled surface $X$ of \cref{lem:ruled-explicit-nef}, keep its nef curvature form $\theta$, its section $D$ at infinity and its K\"ahler form $\kappa=\theta+\pi^*\sigma_0$. For every smooth positive probability form $\mu$,
\[
 \mathcal V_\mu(\theta+\epsilon\kappa)=2\epsilon(1+\epsilon),
 \qquad \mathcal V_\mu(\epsilon\kappa)=2\epsilon^2.
\]
The only positive current of the form $\theta+\ddc u$ is $[D]$. It has minimal singularities, and removing its whole divisor leaves the zero background and changes the common order from one to two.
\end{corollary}
\begin{proof}
\Cref{lem:ruled-explicit-nef} gives the uniqueness of $[D]$ and the nefness of $\theta$, and the calculation in the proof of \cref{prop:nonrecovery-zero-residual-volume} gives $\int_X\theta^2=0$. Put $B=\pi^*\sigma_0$. Then $B^2=0$, and
\[
 \int_X\theta\wedge B=\int_D B=1,
\]
by the Poincar\'e--Lelong formula, Stokes' theorem and the fact that the section $D$ maps isomorphically onto the elliptic base. The closed nef formula following \cref{thm:gaud-volume-nef-limit} now computes both volumes by expanding the squares, for every $\mu$. Every $\theta$-plurisubharmonic function $u$ satisfies $\theta+\ddc u=[D]$, so any two such functions differ by a pluriharmonic function, hence by a constant; in particular each has minimal singularities. Since $\theta$ is itself a curvature form of $\mathcal O_X(D)$, removing the whole divisor leaves the zero background. The two volume polynomials give the orders one and two.
\end{proof}
